\subsection*{The notion of truth in mathematics}

Mathematicians have always been convinced that what they prove is ``true''. It is clear that such a conviction can be only of a sentimental or metaphysical order, and cannot be justified, or even ascribed a meaning which is not tautological, within the domain of mathematics. The history of the concept of truth in mathematics therefore belongs to the history of philosophy and not of mathematics; but the evolution of this concept has had an undeniable influence on the development of mathematics, and for this reason we cannot pass over it in silence.

First of all, a mathematician who is well schooled in philosophy is as rare as a philosopher with a wide knowledge of mathematics. The opinions of mathematicians on philosophical questions, even questions which touch on their own science, are most often second- or third-hand and from sources of doubtful value. But, precisely for this reason, these ``average'' opinions are of as much interest to the historian of mathematics as the original views of thinkers such as Descartes or Leibniz (to cite two who were also mathematicians of the highest order), Plato (who was at least familiar with the mathematics of this time), and Aristotle or Kant (for whom one cannot say as much).

The traditional notion of mathematical truth goes back to the Renaissance. In this conception, there is no great difference between the objects studied by mathematics and those studied by the natural sciences. Both are knowable, and can be apprehended by intuition and by reason; neither intuition nor reason is fallible if properly employed. ``Only an entirely false intellect'', says Pascal, ``could reason wrongly from principles so obvious that it is almost impossible for them to escape notice'' ({[}11{]}, vol.~XII, p.~9). Descartes convinced himself that ``it is only mathematicians who have been able to find proofs, that is to say, certain and evident reasons'' ({[}10{]}, vol.~VI, p.~19), and this (if we accept his account) before he had constructed a metaphysic in which he says, ``this, which a little while ago I took as a rule, namely that the things we conceive very clearly and very distinctly are all true, is assured only because God is or exists and is a perfect being'' ({[}10{]}, vol.~VI, p.~38). Although Leibniz objects that it is not obvious how one is to recognize that an idea is ``clear and distinct'' (*), yet he too considers the axioms to be self-evident and ineluctable consequences of the definitions as soon as the terms are understood (†). Moreover, it should not be forgotten that at this period mathematics embraced many sciences --- sometimes even the art of the engineer which we would not nowadays regard as mathematical. The remarkable success of its applications to ``natural philosophy'', the ``mechanical arts'', and navigation was a principal cause of the confidence it inspired.

From this point of view the axioms were no more open to discussion or doubt than the rules of logical reasoning; at most, it was a matter of personal choice whether one reasoned ``in the manner of the ancients'' or gave free rein to one's intuition. The choice of starting point was also a matter of individual preference. Thus many editions of Euclid's work were published in which the solid logical framework of the Elements was strangely travestied, and ``deductive'' accounts of infinitesimal calculus or rational mechanics were given in which the foundations were singularly badly laid. Spinoza therefore perhaps acted in good faith when he claimed that his Ethics was ``proved in the manner of the geometers'' (more geometrico demonstrata). Although it is hard to find two mathematicians in the seventeenth century in agreement on any single question whatsoever, although the controversies were frequent, interminable, and acrimonious, yet the notion of truth was free from all imputation. ``Since there is only one truth for each thing'', said Descartes, ``whoever finds it knows as much as can be known'' ({[}10{]}, vol.~VI, p.~21).

Although no Greek mathematical text of the great period on these questions has been preserved, it is probable that the point of view of the Greek mathematicians on this subject was much less rigid. It is only by experience that the rules of reasoning have been developed to the point of inspiring complete confidence; before they could be considered as above all discussion they must have passed through many hesitations and tentative forms. Furthermore, it would have been out of keeping with the critical spirit of the Greeks and their taste for discussion and sophistry, if those ``axioms'' which Pascal considered to be self-evident (and which, according to a legend originated by his sister, he himself discovered, with an infallible instinct, as a child) had not been the object of long discussions. In a domain which is not, strictly speaking, that of geometry, the paradoxes of the Eleatic philosophers have preserved for us some traces of such controversies; and Archimedes, when he observes ({[}4{]}, p.~265) that his predecessors have on many occasions made use of the axiom which bears his name, adds that what is proved with the help of this axiom ``has been admitted no less than what is proved without it'', and that it is enough for him that his own results should be admitted on the same footing. Plato, in conformity with his metaphysical views, presented mathematics as a means of access to ``truth in itself'', and the objects with which it deals as existing in their own right in the world of ideas. He characterizes the mathematical method with accuracy in a famous passage of the Republic: ``Those who study geometry and arithmetic\ldots{} assume the existence of odd and even numbers, and three kinds of angles; these things they take as known, and consider that there is no need to justify them either to themselves or to others, because they are self-evident to everyone; and, starting from them, they proceed consistently step by step to the propositions which they set out to examine'' (Book VI, 510c-e). Thus what constitutes a proof is first of all a starting point, which is to some extent arbitrary (although ``self-evident to everyone'') and from which, he says a little later, one does not attempt to go further back; next, a path which passes in order through a sequence of intermediate stations; finally, at each step the consent of the interlocutor guaranteeing the correctness of the reasoning. It should be added that, once the axioms had been stated, no new appeal to intuition was in principle admitted; Proclus, quoting Geminus, recalls that ``we have taught even the pioneers in this science to take no account of merely plausible conclusions when concerned with reasoning that is to form a part of our geometric doctrine'' ({[}3{]}, vol.~I, p.~203).

Thus the rules of mathematical reasoning were developed by experience and the fire of criticism. If it is true, as has been plausibly argued (*), that Book VIII of Euclid has preserved for us a part of the arithmetic of Archytas, then it is not surprising to find there the somewhat pedantic stiffness of reasoning which never fails to appear in any mathematical school which discovers or thinks it has discovered ``rigour''. But once these rules of reasoning had entered into the common practice of mathematicians, it does not appear that they were ever subjected to doubt until a very recent epoch; although Aristotle and the Stoics deduced certain of these rules from certain others by schemes of reasoning, the primitive rules were always admitted as self-evident. Likewise, having gone back as far as the ``hypotheses'', ``axioms'', and ``postulates'' which seemed to them to provide a solid foundation for the science of their period (such, for example, as must have appeared in the first Elements, which tradition ascribes to Hippocrates of Chio, about 450 B.C.), the Greek mathematicians of the classical period seem to have devoted their efforts to the discovery of new results rather than to a critique of the foundations, which at this period could only have been sterile; and, metaphysical preoccupations apart, the text from Plato quoted earlier bears witness to this general agreement among mathematicians concerning the bases of their science. Furthermore, the Greek mathematicians seem not to have believed that the ``primary notions'' which served for them as a starting point --- straight line, surface, ratio of magnitudes --- were capable of further elucidation. If they gave ``definitions'', it was visibly to clear their consciences and without any illusions about the validity of these definitions. It goes without saying that, by contrast, the Greek mathematicians and philosophers had perfectly clear ideas on definitions other than those of the ``primary notions'' (often called ``nominal'' definitions, and playing the same role as our ``abbreviating symbols''). In this connection there intervenes explicitly, without doubt for the first time, the question of ``existence'' in mathematics. Aristotle did not omit to remark that a definition does not imply the existence of the thing defined, and that either a postulate or a proof of existence is necessary. Undoubtedly his observation was derived from the practice of mathematicians. In every case Euclid is careful to postulate the existence of the circle, to prove that of the equilateral triangle, parallel lines, the square, etc. as he finds it necessary to introduce them into his arguments ({[}3{]}, Book I); these proofs are ``constructions'' --- in other words, he exhibits, with the aid of his axioms, mathematical objects which he then proves satisfy the definitions in question.

Thus we see Greek mathematics of the classical period come to a sort of empirical certainty (whatever may have been the metaphysical basis, according to this or that philosopher). If the rules of mathematical reasoning are considered to be above question, then the success of Greek science, and the feeling that critical revision would have been inopportune, are largely due to the confidence inspired by the axioms: a confidence of the same order as that accorded in the last century to the principles of theoretical physics. It is what is suggested by the proverb ``nihil est in intellectu quod non prius fuerit in sensu'' with which Descartes rightly took issue as not providing a firm enough basis for what he proposed to achieve by the use of reason.

It was not until the beginning of the nineteenth century that mathematicians retreated from the arrogant position of Descartes (not to speak of Kant or of Hegel; the latter was a little behind the time, as was to be expected, in relation to the science of his age (*)) to a position as flexible as that of the Greeks. The first blow to the classical conceptions was the construction of hyperbolic non-Euclidean geometry by Gauss, Lobatschevsky, and Bolyai at the beginning of the century. We shall not undertake here a detailed account of the genesis of this discovery, the culmination of numerous fruitless attempts to prove the parallel postulate. At the time, its effect on the principles of mathematics was perhaps not so deep as is sometimes asserted. It simply demanded the abandonment of the claims of the preceding century to the ``absolute truth'' of Euclidean geometry and, a fortiori, of the Leibnizian view that the definitions imply the axioms; the latter no longer appear at all as ``self-evident'' but rather as hypotheses which are to be judged according to whether they are adapted to the mathematical representation of the physical world. Gauss and Lobatschevsky believed that the debate between the various possible geometries could be settled by experiment ({[}15{]}, p.~76). This was also the point of view of Riemann, the aim of whose famous inaugural lecture ``On the hypotheses which lie at the foundations of geometry'' was to provide a general mathematical framework for various natural phenomena. ``The question remains to be settled'', he said, ``as to how far and to what extent these hypotheses will be confirmed by experiment'' ({[}19{]}, p.~284). But clearly this is a problem which no longer has anything to do with mathematics; and none of the authors just mentioned seems to have had any doubts that, even if a ``geometry'' did not correspond to experimental reality, its theorems were any less ``mathematically true'' (*).

However, if this is so, such a conviction can no longer be attributed to an unlimited confidence in classical ``geometrical intuition''. The description which Riemann sought to give of ``manifolds n times extended'', the object of his work, relies on ``intuitive'' considerations (†) only to the point of justifying the introduction of ``local coordinates''; from then on, he seems to have felt that he was on firm ground, namely that of Analysis. But the latter is ultimately based on the concept of real number, which up to that time had remained on a very intuitive basis; and progress in the theory of functions was leading to most disturbing results in this respect. With the work of Riemann himself on integration, and even more with the examples of curves with no tangents, constructed by Bolzano and Weierstrass, the whole pathology of mathematics was beginning to emerge. A century later, we have seen so many monsters of this sort that we have become a little blasé. But the effect produced on the majority of nineteenth century mathematicians ranged from disgust to consternation. ``How'', asked H. Poincaré, ``can intuition deceive us at this point?'' ({[}33 b{]}, p.~19); and Hermite (not without a turn of humour, which not all the commentators on this celebrated phrase seem to have perceived) declared that he ``turned away with fright and horror from this lamentable plague of continuous functions with no derivatives'' ({[}27{]}, vol.~II, p.~318). The most serious feature of the malaise was that it was no longer possible to attribute these phenomena, so contrary to common sense, to badly elucidated notions, as in the days of the ``indivisibles'', because they arose after the reforms of Bolzano, Abel, and Cauchy, which had succeeded in founding the notion of limit as rigorously as the theory of proportions. They were therefore to be ascribed to the gross and incomplete character of our geometrical intuition, which since that time has been discredited as a method of proof.

This conclusion inevitably affected classical mathematics, beginning with geometry. However much respect was paid to the axiomatic construction of Euclid, the imperfections in it had been noticed, from antiquity onward. The parallel postulate had been the object of the largest number of criticisms and attempts at proof; but Euclid's successors and commentators had also sought to prove other postulates (notably that of the equality of right angles) or had recognized the inadequacy of certain definitions, such as those of a straight line or a plane. In the 16th century Clavius, an editor of the Elements, noticed the absence of a postulate guaranteeing the existence of the fourth proportional; Leibniz, for his part, remarked that Euclid makes use of geometrical intuition without mentioning it explicitly, for example when he assumes (Elements, Book I, Proposition 1) that two circles, each of which passes though the centre of the other, have a common point ({[}12 b{]}, vol.~VII, p.~166). Gauss (who did not deny himself the use of such topological considerations) drew attention to the part played in Euclid's constructions by the notion of a point (or a line) situated ``between'' two others, a notion which nevertheless had not been defined ({[}14{]}, vol.~VIII, p.~222). Finally, the use of displacements --- notably in the case of ``congruent triangles'' --- long accepted as standing to reason (*), soon appeared to the critics of the nineteenth century as a concept which rested on unformulated axioms. Thus, in the period from 1860 to 1885, many partial revisions of the foundations of geometry appeared (Helmholtz, Méray, Houël) with the purpose of filling some of these gaps. But it was not until M. Pasch {[}28{]} that the abandonment of all appeal to intuition was clearly formulated as a programme and carried out with perfect rigour. The success of his enterprise soon produced numerous emulators who, principally between 1890 and 1910, gave varied presentations of the axioms of Euclidean geometry. The most famous of these works were Peano's, written in his symbolic language {[}29 b{]}, and above all Hilbert's Grundlagen der Geometrie, which appeared in 1899, and by the lucidity and profundity of its exposition immediately and justly became the model of modern axiomatics, to the extent that its predecessors were forgotten. Hilbert was not content with giving a complete system of axioms for Euclidean geometry, but classed them in various groups according to their nature, and set about determining the exact range of each group of axioms, not only by developing the logical consequences of each group in isolation, but also by discussing the various ``geometries'' obtained by suppressing or modifying certain axioms (geometries among which those of Lobatschevsky and Riemann appear as special cases (*)). In this way he brought clearly into view, in a domain which until that time had been considered as one of the closest to external reality, the liberty which a mathematician has in the choice of his postulates. In spite of the disarray caused among some philosophers by these ``metageometries'' with their weird properties, the thesis of the Grundlagen was rapidly and almost unanimously adopted by mathematicians. Thus H. Poincaré, who could hardly be suspected of partiality towards formalism, asserted in 1902 that the axioms of geometry are conventions for which the notion of ``truth'', in the everyday sense of the word, has no meaning ({[}33 a{]}, pp.~66-67). ``Mathematical truth'' therefore consists entirely in logical deduction from premises posed arbitrarily as axioms. As we shall see later, the validity of the rules of reasoning according to which the deductions operate was soon to be called in question, leading thus to a complete overhauling of the conceptions at the base of mathematics.

